\section{Square-lattice arithmetic, incomplete gamma and uniform Fourier bounds}
\label{sec:squarelattice}
The previous generic Gaussian--hyperbolic counterexamples
show why shape restrictions cannot by themselves locate
Fourier zeros. Here the exact square-lattice weights
$q_m=\pi m^2$ restore genuinely arithmetic information.
We distinguish three facts: a classical incomplete-gamma
resummation, a quantitative \emph{uniform in real frequency}
truncation theorem, and a modular cancellation mechanism
that each fixed truncation necessarily loses.

\paragraph{An exact incomplete-gamma expansion.}
Let $q_m=\pi m^2$ and, on $r\ge0$, put
\begin{equation}
\phi_m(r)=q_m(2q_me^{4r}-3)e^{5r-q_me^{4r}}.
\label{eq:arithmeticmode}
\end{equation}
Write $\Phi(r)=\sum_{m\ge1}\phi_m(r)$,
$F(z)=\int_{\mathbb R}\Phi(r)e^{izr}\,dr$,
and define the even positive-axis truncation
\[
F_M(z)=2\int_0^\infty
\left(\sum_{m=1}^{M}\phi_m(r)\right)
\cos(zr)\,dr.
\]
Then $F(z)=\xi(1/2+iz/2)/4$ in our fixed normalization.
For $a_\pm=5/4\pm iz/4$, define
\[
H_m(a)=q_m^{1-a}
\bigl(2\Gamma(a+1,q_m)-3\Gamma(a,q_m)\bigr).
\]
With $\Gamma(a,q)=\int_q^\infty t^{a-1}e^{-t}\,dt$,
the substitution $t=e^{4r}$ gives the
\emph{exact analytic identity}
\begin{equation}
F_M(z)=\frac14\sum_{m=1}^{M}
\bigl(H_m(a_+)+H_m(a_-)\bigr).
\label{eq:arithmeticgamma}
\end{equation}
On the real axis this becomes
$F_M(x)=\frac12\Re\sum_{m\le M}H_m(5/4+ix/4)$.
The series converges locally uniformly in $z$;
incomplete-gamma representations of $\xi$ are established
prior art \cite{DF,Haglund,Polson}, so
\eqref{eq:arithmeticgamma} is used as an analytic
tool rather than claimed as a new representation.

\begin{theorem}[Uniform real-frequency derivative and Laguerre
truncation]
\label{thm:uniformFourier}
For every integer $M\ge1$, derivative order $k\ge0$,
and real $x$, let $q=q_{M+1}=\pi(M+1)^2$. Then
\begin{equation}
\boxed{\left|F^{(k)}(x)-F_M^{(k)}(x)\right|
\le \frac{3k!}{4^k}q e^{-q}.}
\label{eq:unifderiv}
\end{equation}
The constant
\[
C=e^{-\pi}
\left(\pi+3+\frac6\pi+\frac6{\pi^2}\right)
+12\pi e^{-4\pi}
\]
also satisfies
$|F^{(k)}(x)|,|F_M^{(k)}(x)|
\le Ck!/4^k$ for all $k,M,x$.
Consequently, for every $n\ge0$ and $x\in\mathbb R$,
\begin{equation}
\boxed{\left|L_n[F](x)-L_n[F_M](x)\right|
\le
\frac{6(2n+1)C}{4^{2n}}
q_{M+1}e^{-q_{M+1}}.}
\label{eq:uniformlag}
\end{equation}
\end{theorem}
\begin{proof}
The theta summands are positive for $r\ge0$, since
$q_m>3$ and $2q_me^{4r}-3>0$.
For $k\ge0$, Taylor's exponential inequality gives
$r^k\le k!e^{4r}/4^k$.
Differentiating the Fourier transform under the
integral and taking absolute values bounds the
omitted tail by
\begin{align*}
&\frac{2k!}{4^k}
\sum_{m>M}\int_0^\infty
2q_m^2e^{13r-q_me^{4r}}\,dr\\
&\quad=\frac{k!}{4^k}
\sum_{m>M}q_m^2
\int_1^\infty t^{9/4}e^{-q_m t}\,dt.
\end{align*}
As $t\ge1$, $t^{9/4}\le t^3$.
Integrating the integer-degree upper bound yields
\begin{equation}
\frac{k!}{4^k}\sum_{m>M}e^{-q_m}
\left(q_m+3+\frac6{q_m}+\frac6{q_m^2}\right).
\label{eq:tailterm}
\end{equation}
For $m\ge2$, $q_m>12$ and elementary
rational comparison gives
\[
q_m+3+6/q_m+6/q_m^2\le 3q_m/2.
\]
Moreover
\[
\frac{q_{m+1}e^{-q_{m+1}}}
{q_me^{-q_m}}
=\left(\frac{m+1}{m}\right)^2
e^{-\pi(2m+1)}
<4e^{-9}<\tfrac12.
\]
Summing this geometric majorant proves
\eqref{eq:unifderiv}. Applying the same
estimate to the initial $m=1$ term and the
tail $m\ge2$ proves the stated $C$ bound
for both $F$ and $F_M$.

Write $E_k=F^{(k)}-F_M^{(k)}$.
Expanding the difference of the products in
the definition of $L_n$ yields
\[
E_kF^{(2n-k)}
+F_M^{(k)}E_{2n-k}.
\]
The derivative bounds, together with
\[
\binom{2n}{k}k!(2n-k)!=(2n)!,
\]
turn each of its $2n+1$ summands into at most
$2C(3q e^{-q})/4^{2n}$ after factorial
normalization. This proves \eqref{eq:uniformlag}.
\end{proof}

\section{An infinite modular odd-jet cancellation law}
\label{sec:modularoddjets}
Let $\psi(t)=\sum_{m\ge1}e^{-\pi m^2t}$,
and put $h(r)=e^r\psi(e^{4r})$.
The Jacobi theta transformation gives
\begin{equation}
h(r)-h(-r)=-\sinh r.
\label{eq:modularh}
\end{equation}
Direct differentiation of one summand shows
\[
(D_r^2-1)h(r)=8\Phi(r).
\]
As $(D_r^2-1)\sinh r=0$,
the functional equation \eqref{eq:modularh}
proves $\Phi(-r)=\Phi(r)$ without assuming
term-by-term evenness.

For $P_0(q)=2q-3$ and
$P_{j+1}=(5-4q)P_j+4qP_j'$,
termwise differentiation gives
\[
\phi_m^{(j)}(0)=q_me^{-q_m}P_j(q_m).
\]
Evenness therefore establishes the
\emph{infinite exact arithmetic identities}
\begin{equation}
\boxed{\sum_{m\ge1}
q_me^{-q_m}P_{2j+1}(q_m)=0,
\quad j=0,1,2,\ldots.}
\label{eq:oddjetidentity}
\end{equation}
At $j=0$ this becomes
\begin{equation}
\boxed{
\sum_{m\ge1}q_me^{-q_m}
(-8q_m^2+30q_m-15)=0.}
\label{eq:firstmodularcancellation}
\end{equation}
These identities are consequences of
the classical theta modular transformation,
not new modular equations.

\begin{theorem}[Every fixed theta cutoff produces
an eventual false first-Laguerre signal]
\label{thm:falsecutoff}
For every finite $M\ge1$, define
$f_M(r)=\sum_{m=1}^M\phi_m(r)$ on $r\ge0$,
and $F_M(x)=2\int_0^\infty f_M(r)\cos(xr)\,dr$.
Then
\[
D_M:=f_M'(0)>0
\]
and, as real $|x|\to\infty$,
\begin{align}
F_M(x)&=-2D_Mx^{-2}+O(x^{-4}),\nonumber\\
F_M'(x)&=4D_Mx^{-3}+O(x^{-5}),\nonumber\\
F_M''(x)&=-12D_Mx^{-4}+O(x^{-6}).
\end{align}
Consequently
\begin{equation}
\boxed{
L_1[F_M](x)
=-8D_M^2x^{-6}+O(x^{-8})<0
}
\label{eq:spurious}
\end{equation}
for all sufficiently large real $|x|$.
In contrast, the complete modular kernel $\Phi$
has vanishing odd endpoint derivatives of
\emph{every} order and its Fourier transform
decays faster than any inverse power
on the real axis.
\end{theorem}
\begin{proof}
The polynomial
$P_1(q)=-8q^2+30q-15$ is strictly negative
for $q\ge12$: writing $q=12+t$, $t\ge0$,
\[
-P_1(q)=8t^2+162t+807>0.
\]
Every omitted lattice index $m>M\ge1$
has $q_m=\pi m^2\ge4\pi>12$, so
the omitted terms in
\eqref{eq:firstmodularcancellation}
are strictly negative. The full sum is
zero; therefore $D_M$ is strictly positive.
All derivatives of the finite smooth
positive-axis sum decay integrably at
infinity. Repeated integration by parts
in the cosine integral produces the
three stated endpoint asymptotics.
Substitution into
$L_1=(F')^2-FF''$ gives \eqref{eq:spurious}.
For the complete kernel, the modular parity
\eqref{eq:oddjetidentity} eliminates
\emph{all} boundary terms; integrating
by parts any fixed number of times
establishes superalgebraic Fourier decay.
\end{proof}

\paragraph{Numerical replay and its limits.}
The following $55$-digit mpmath evaluations
at $x=39.5$ independently use the original
Fourier integral and incomplete-gamma
representation:
\begin{center}\small
\begin{tabular}{cc}
\hline
Mode cutoff $M$&$L_1[F_M](39.5)$\\
\hline
1 & $-5.74158212233354\,10^{-11}$\\
2 & $+3.24498665006940\,10^{-11}$\\
3 & $+3.24499162628817\,10^{-11}$\\
\hline
\end{tabular}
\end{center}
They illustrate a sign reversal between
the first and second square-lattice modes;
these sampled values are \emph{numerical
evidence}, not directed-rounding sign
certificates. The identities
\eqref{eq:firstmodularcancellation} and
\eqref{eq:spurious}, by contrast, are
analytic theorems.

\paragraph{Interpretation.}
The quantitative bounds
\eqref{eq:unifderiv}--\eqref{eq:uniformlag}
converge uniformly over \emph{all real}
Fourier frequencies for each fixed
differentiation or Laguerre order.
Nevertheless every finite truncation
eventually has negative $L_1$.
This is not a contradiction: the
genuine transform and its Laguerre
expressions approach zero at large
frequencies, so an absolute uniform
error bound alone cannot guarantee
their \emph{sign} globally.
A successful RH approach must preserve
or control the exact modular cancellation,
not rely on a fixed-mode sign extrapolation.

\paragraph{Verification and prior art.}
The scripts
\texttt{certificates/xi\_square\_lattice\_uniform\_tail.py},
\texttt{certificates/xi\_modular\_odd\_jet\_exact.py},
and
\texttt{certificates/xi\_arithmetic\_incomplete\_gamma\_replay.py}
provide exact rational inequality checks
and independent high-precision
replays, with each status distinguished.
The Lean modules
\texttt{XiKernel/SquareLatticeFourierTail.lean}
and \texttt{XiKernel/ModularOddJet.lean}
certify the finite polynomial gates;
the theta functional equation and all
Fourier asymptotics remain analytical
paper proofs, not end-to-end Lean theorems.
Classical theta transformation,
incomplete-gamma representations,
and truncated-\(\Xi\) investigations
have substantial prior art
\cite{DF,Haglund,Polson};
originality of this precise collection
of bounds and truncation consequences
has not been established.
